% =============================================================================
%  CAPÍTULO 4 — RESULTADOS E DISCUSSÃO
%  TCC: Dollar Bars vs. Time Bars — Modelos GARCH para BTC/USDT
% =============================================================================

\chapter{Resultados e Discussão}
\label{cap:resultados}

% -----------------------------------------------------------------------------
\section{Propriedades Estatísticas Comparativas das Séries}
\label{sec:prop_estatisticas}

A análise exploratória das séries de retornos logarítmicos revela diferenças estruturais
substanciais entre as duas formas de amostragem, conforme sintetizado na Tabela~\ref{tab:estat_descritiva}.
Essas diferenças não são ruído amostral: elas decorrem diretamente do mecanismo pelo qual
cada tipo de barra é construída.

O indicador mais expressivo é a curtose. As time bars de 5 minutos apresentam curtose de Fisher
de 75,37, ao passo que as dollar bars registram apenas 3,09 — uma razão de aproximadamente
25 vezes. Para compreender essa disparidade, é preciso considerar o que ocorre durante
choques de volatilidade no mercado de Bitcoin. Ao longo de um episódio de alta atividade
financeira, dezenas de transações são comprimidas em uma única barra temporal, produzindo
retornos extremamente grandes que inflariam artificialmente as caudas da distribuição.
A dollar bar, ao encerrar uma barra sempre que o volume financeiro transacionado atinge o
limiar $\bar{D}$, redistribui esse mesmo choque ao longo de múltiplas barras, diluindo o
extremo. Esse mecanismo é precisamente a sub-amostragem em períodos de alta atividade
descrita por \citet{Prado2018} na discussão sobre barras de informação (seção 2.1.1).

A assimetria reforça o mesmo argumento. Enquanto as dollar bars exibem assimetria de apenas
0,05 — praticamente simétrica —, as time bars apresentam assimetria de $-1{,}01$. O sinal
negativo indica que quedas abruptas são capturadas de forma concentrada nessa amostragem:
grandes retornos negativos, correspondentes a correções rápidas, tendem a ocorrer em janelas
de poucos minutos e, por isso, se concentram em uma única barra temporal. Além da implicação
distribucional, a simetria das dollar bars facilita a especificação do modelo GARCH, que,
em sua formulação padrão com inovações gaussianas, assume distribuição condicional simétrica.

Os testes formais de adequação distribucional, apresentados na Tabela~\ref{tab:testes_formais},
amplificam esse quadro. A estatística de Jarque-Bera das time bars é superior a 49 milhões,
contra aproximadamente 83 mil para as dollar bars — uma diferença de mais de 500 vezes.
Esse resultado não representa apenas rejeição da normalidade, que ambas as séries exibem:
indica que as caudas das time bars pertencem a uma ordem de magnitude distinta. O teste
de Ljung-Box nos retornos em nível rejeita a hipótese de ausência de autocorrelação
para as duas séries, mas com estatística consideravelmente maior para as time bars
(155,03 versus 50,06). Contudo, conforme discutido na seção seguinte, essa autocorrelação
é de magnitude modesta para o volume de observações disponível e não justifica componentes
autorregressivos ou de médias móveis no modelo final. Ambas as séries são estacionárias:
o teste Dickey-Fuller Aumentado rejeita a hipótese de raiz unitária, e o teste KPSS não
rejeita estacionariedade em nenhuma das duas amostras, validando o uso de $d = 0$ na
especificação ARMA.

Em conjunto, os resultados exploratórios corroboram a hipótese de \citet{Prado2018}:
dollar bars produzem séries com propriedades distribucionais superiores em termos de
curtose, assimetria e dependência linear residual. Cabe ressalvar, contudo, que propriedades
descritivas melhores não garantem, por si só, maior poder preditivo fora da amostra.
Essa questão é investigada nas seções seguintes.

% -----------------------------------------------------------------------------
\section{Seleção de Modelos via CPCV}
\label{sec:selecao_modelos}

O procedimento de seleção de hiperparâmetros via \textit{Combinatorial Purged Cross-Validation}
(CPCV) com $N = 10$ grupos, $K = 5$ grupos de teste e embargo de 240 barras gerou
$\binom{10}{5} = 252$ folds de avaliação, totalizando 36.288 ajustes MLE por série
(252 folds $\times$ 144 combinações de parâmetros no grid
$\text{ARMA}(0..5,\,0..5) \times \text{GARCH}(1..2,\,1..2)$).

Em ambas as séries, o componente ARMA selecionado foi ARMA$(0,0)$, ou seja, média constante
sem termos autorregressivos nem de médias móveis. Esse resultado é consistente com a
hipótese de eficiência de mercado na forma fraca: os retornos do par BTC/USDT, a despeito
de sua curtose elevada, não apresentam estrutura linear preditível suficiente para justificar
a inclusão de termos AR ou MA. É importante observar que esse resultado não contradiz a
rejeição do Ljung-Box reportada na seção anterior. As autocorrelações detectadas são de
magnitude prática pequena — estatísticas de 50 e 155 para séries com cerca de 200 mil
observações representam efeitos de baixa amplitude — e concentram-se, sobretudo, nos
retornos ao quadrado, não nos níveis; o componente GARCH é o mecanismo mais adequado
para capturá-las.

O resultado mais informativo da seleção refere-se ao componente GARCH. Para as dollar bars,
o modelo GARCH$(1,2)$ — com um termo de inovação ao quadrado $(\alpha_1)$ e dois termos
de variância defasada $(\beta_1, \beta_2)$ — foi escolhido em 233 dos 252 folds,
correspondendo a 92,5\% de frequência de seleção. Para as time bars, o modelo GARCH$(1,1)$
canônico foi selecionado na totalidade dos 252 folds (100,0\%). O AIC médio das time bars
ao longo dos folds foi de $-859.083$, contra $-835.475$ para as dollar bars, indicando
melhor ajuste \textit{in-sample} para a série temporal.

A necessidade de um termo $\beta_2$ adicional nas dollar bars tem uma interpretação
econômica plausível. Por construção, as dollar bars possuem espaçamento temporal irregular:
são geradas em maior número durante períodos de alta atividade financeira, quando a
volatilidade do mercado é tipicamente mais elevada. Esse mecanismo faz com que os choques
de volatilidade sejam distribuídos ao longo de mais barras, tornando o decaimento da
variância condicional mais gradual e demandando uma especificação ligeiramente mais rica
para capturá-lo adequadamente. As time bars, ao ignorar esse padrão de irregularidade
temporal, comprimem os choques em barras únicas e produzem uma dinâmica de volatilidade
mais diretamente capturada pelo GARCH$(1,1)$ canônico.

A estabilidade das seleções merece atenção. O fato de o GARCH$(1,1)$ ter sido eleito
em 100\% dos folds para as time bars indica que a dinâmica de volatilidade dessa série
é bem descrita pela especificação padrão da literatura, com pouca ambiguidade quanto à
ordem. Para as dollar bars, a frequência de 92,5\% tampouco é marginal: denota robustez
do GARCH$(1,2)$, não um artefato de um fold específico. Em síntese, o CPCV entregou
seleções estáveis para ambas as séries, com a diferença de especificação entre elas
refletindo uma propriedade genuína das dinâmicas de volatilidade subjacentes.

Cabe destacar que ajuste \textit{in-sample} superior não implica, necessariamente,
melhor desempenho preditivo fora da amostra — distinção fundamental em qualquer
avaliação de modelos de previsão \citep{Diebold2015}. Essa questão é endereçada
diretamente na seção seguinte.

% -----------------------------------------------------------------------------
\section{Desempenho Preditivo Fora da Amostra}
\label{sec:desempenho_previsao}

Esta seção constitui o núcleo analítico do trabalho, respondendo diretamente à pergunta
de pesquisa com base nas métricas calculadas sobre os conjuntos de teste — 171 dias para
as dollar bars (2025-07-14 a 2025-12-31) e 147 dias para as time bars
(2025-08-07 a 2025-12-31) — e comparadas contra o estimador TSRV de volatilidade
realizada como benchmark.

\subsection{Métricas Pontuais}
\label{subsec:metricas_pontuais}

As métricas de avaliação pontual, apresentadas na Tabela~\ref{tab:metricas_comparacao},
revelam dominância das dollar bars nas três funções de perda avaliadas. O Erro Quadrático
Médio (MSE) das dollar bars é de $1{,}08 \times 10^{-7}$, contra $1{,}65 \times 10^{-7}$
para as time bars — uma redução de aproximadamente 34\%. O Erro Absoluto Médio (MAE)
exibe comportamento análogo: $2{,}72 \times 10^{-4}$ versus $3{,}59 \times 10^{-4}$,
representando redução de aproximadamente 24\%. A métrica QLIKE, cujos valores menores
(em módulo, mais negativos) indicam melhor desempenho, registra $-7{,}026$ para as dollar
bars e $-6{,}776$ para as time bars, uma diferença de 0,25 em magnitude absoluta.

As três métricas pertencem à classe de funções de perda robustas de \citet{Patton2011},
cuja propriedade fundamental é preservar o ranking entre previsões de volatilidade mesmo
quando o benchmark utilizado como variável dependente é um proxy ruidoso da volatilidade
latente verdadeira. Essa propriedade é particularmente relevante no presente contexto,
pois o estimador TSRV, ainda que consistente, incorpora ruído de microestrutura e
incerteza de estimação que o afastam da variância quadrática integrada teórica. O fato
de as três métricas robustas apontarem unanimemente na mesma direção reforça a
confiabilidade do resultado.

\subsection{Regressão de Mincer-Zarnowitz}
\label{subsec:mz}

A regressão de Mincer-Zarnowitz \citep{MincerZarnowitz1969}, que regride a volatilidade
realizada diária $\hat{\sigma}^2_{\text{RV},t}$ sobre a previsão do modelo
$\hat{h}_t$, fornece um diagnóstico adicional sobre a qualidade e o viés das previsões.
Para as dollar bars, o coeficiente de determinação $R^2$ foi de 0,655 e o intercepto
estimado de $\hat{\alpha} = -1{,}08 \times 10^{-4}$ (p-valor = 0,0022). Para as time bars,
$R^2 = 0{,}493$ e $\hat{\alpha} = -2{,}54 \times 10^{-4}$ (p-valor = 0,0001).

O modelo baseado em dollar bars explica, portanto, 65,5\% da variação da volatilidade
realizada diária, contra 49,3\% para as time bars. Essa diferença de aproximadamente
16 pontos percentuais no $R^2$ é economicamente relevante: para aplicações de gestão
de risco, como o cálculo de VaR (\textit{Value at Risk}) ou a otimização de carteiras
com restrições de variância, modelos que explicam maior parcela da variação da volatilidade
realizada tendem a produzir limites de risco mais calibrados \citep{Andersen2003}.

Ambos os interceptos são negativos e estatisticamente significativos, indicando leve
subestimação sistemática da volatilidade realizada por parte de ambos os modelos.
Esse resultado é típico de modelos GARCH paramétricos estimados com inovações gaussianas
quando comparados a estimadores não-paramétricos de alta frequência: o estimador TSRV
captura componentes de variação intrajornada que o GARCH, por ser estimado sobre retornos
diários ou intradiários agregados, não incorpora diretamente. O coeficiente angular $\hat{\beta}$
foi de 0,858 para as dollar bars (p-valor $\approx 0$) e de 1,032 para as time bars
(p-valor $\approx 0$), sugerindo que o modelo de dollar bars tende a subestimar
a resposta da volatilidade realizada à previsão, enquanto o de time bars apresenta
resposta ligeiramente acima da unidade, indicando leve sobre-reação relativa.

\subsection{Teste de Diebold-Mariano}
\label{subsec:dm}

A comparação formal de capacidade preditiva foi realizada por meio do teste de
Diebold-Mariano \citep{DieboldMariano1995}. A estatística DM calculada sobre o subconjunto
de 147 dias em que ambos os modelos produzem previsões sobrepostas foi de $-1{,}9578$,
com p-valor bilateral de 0,0503.

A interpretação desse resultado requer cuidado analítico. Em primeiro lugar, o sinal
negativo da estatística DM confirma que o modelo de dollar bars apresenta perdas menores
em média: a direção da diferença é inequívoca e coerente com todas as demais métricas.
Em segundo lugar, o p-valor de 0,0503 situa-se imediatamente acima do limiar convencional
de 5\%, o que, em sentido estrito, não permite rejeitar a hipótese nula de igualdade
preditiva ao nível de 5\%. Contudo, uma avaliação intelectualmente honesta desse resultado
exige considerar o poder do teste nas condições da amostra disponível.

O teste de Diebold-Mariano é aplicado sobre os 147 dias de sobreposição entre os conjuntos
de teste, não sobre toda a amostra de teste. Com essa quantidade de observações, o poder
estatístico do teste é moderado: pequenas diferenças de desempenho, mesmo que economicamente
relevantes, podem não atingir significância ao nível convencional. Uma amostra de dias de
previsão mais extensa — possível com um horizonte de dados mais longo — provavelmente
produziria rejeição, dada a consistência da vantagem das dollar bars em todas as métricas.
Além disso, conforme argumenta \citet{Diebold2015}, o limiar de 5\% é uma convenção, não
uma fronteira ontológica entre evidência e ausência de evidência: um p-valor de 0,0503
constitui evidência fraca \textit{contra} a hipótese nula de igualdade, não evidência
\textit{a favor} dela. Em termos práticos, uma redução de 34\% no MSE e de 24\% no MAE
é economicamente relevante para um gestor de risco, independentemente do p-valor pontual.

\subsection{Intervalos de Confiança Bootstrap}
\label{subsec:bootstrap}

Os intervalos de confiança de 95\% obtidos por bootstrap com $n_{\text{boot}} = 2000$
reamostras, apresentados na Tabela~\ref{tab:ic_bootstrap}, complementam e, em parte,
superam a conclusão do teste DM ao fornecer incerteza marginal para cada modelo
separadamente.

Para as métricas MSE e MAE, os intervalos se sobrepõem parcialmente entre os dois modelos.
O IC 95\% do MSE das dollar bars é $[7{,}99 \times 10^{-8},\ 1{,}44 \times 10^{-7}]$,
enquanto o das time bars é $[1{,}35 \times 10^{-7},\ 2{,}00 \times 10^{-7}]$. A sobreposição,
embora modesta, é consistente com a não-rejeição pelo teste DM. Para o MAE, situação análoga:
dollar bars em $[2{,}45 \times 10^{-4},\ 3{,}00 \times 10^{-4}]$ e time bars em
$[3{,}29 \times 10^{-4},\ 3{,}90 \times 10^{-4}]$, com sobreposição pequena mas presente.

O resultado mais robusto do presente trabalho emerge da análise dos intervalos de
confiança para a métrica QLIKE. O IC 95\% das dollar bars é $[-7{,}156,\ -6{,}894]$,
e o das time bars é $[-6{,}877,\ -6{,}668]$. Esses intervalos \textbf{não se sobrepõem},
o que significa que, com 95\% de confiança, a métrica QLIKE das dollar bars é menor do
que a das time bars. A relevância desse achado é amplificada pelo fato de o QLIKE ser,
entre as três funções de perda avaliadas, a mais adequada do ponto de vista teórico para
comparar previsões de volatilidade: conforme demonstrado por \citet{Patton2011}, é a
única função de perda da família que é robusta ao uso de proxies ruidosos e que ao mesmo
tempo possui propriedades de elicitabilidade para a variância condicional.

Em síntese, o conjunto de evidências — estatística DM negativa na fronteira de 5\% de
significância, ICs de QLIKE não sobrepostos a 95\%, e dominância simultânea em MSE, MAE,
QLIKE e $R^2$ MZ — constitui evidência consistente e coerente em favor das dollar bars
como insumo para modelos GARCH de previsão de volatilidade do BTC/USDT. A evidência não
é conclusiva ao nível convencional de significância bilateral do DM, mas é robusta em
direção e coerente entre múltiplas perspectivas de avaliação.

% -----------------------------------------------------------------------------
\section{Limitações e Ressalvas}
\label{sec:limitacoes}

A interpretação dos resultados deve ser contextualizada por um conjunto de limitações
metodológicas e de escopo que delimitam o alcance das conclusões.

\paragraph{Ativo único.}
Todos os resultados são obtidos para o par BTC/USDT negociado na Binance.
Criptomoedas operam 24 horas por dia, sete dias por semana, e possuem microestrutura
distinta de mercados regulados: ausência de mecanismos de circuit breaker, inexistência
de leilão de abertura e fechamento, e liquidez fragmentada entre múltiplas exchanges.
Essas características tornam o BTC/USDT um caso de análise com dinâmicas próprias,
e a generalização dos resultados para ações, câmbio ou commodities demanda estudos
adicionais com esses ativos.

\paragraph{Período amostral.}
O horizonte de dois anos (2024–2025) é suficiente para estimação e validação dos modelos,
mas pode não capturar toda a variação de regimes de volatilidade relevantes para o Bitcoin.
O período abrange fases de bull e bear market, mas parâmetros estimados sobre esse intervalo
podem não ser estáveis em regimes futuros com características distintas das observadas.

\paragraph{Calibração do limiar $\bar{D}$.}
O limiar das dollar bars foi calibrado para produzir o mesmo número médio de barras por dia
que as time bars de 5 minutos, tornando as comparações controladas por volume de observação.
Essa escolha é metodologicamente razoável, mas arbitrária: resultados poderiam diferir com
limiares maiores ou menores, que gerariam séries com granularidade distinta e, possivelmente,
propriedades distribucional diferentes.

\paragraph{Distribuição dos erros.}
Ambos os modelos foram estimados com inovações gaussianas. Diante do excesso de curtose
documentado — especialmente expressivo nas time bars — distribuições de cauda pesada,
como a t-Student ou a Distribuição de Erro Generalizado (GED), poderiam melhorar o ajuste
e potencialmente alterar a magnitude e o sentido das diferenças entre os modelos.

% -----------------------------------------------------------------------------
\section{Síntese do Capítulo}
\label{sec:sintese_cap4}

Em conjunto, os resultados indicam que dollar bars produzem séries de retornos com
propriedades estatísticas superiores e modelos GARCH com maior poder preditivo fora
da amostra para o BTC/USDT no período 2024–2025. A evidência é consistente em direção
— dominância em MSE, MAE, QLIKE e $R^2$ MZ — e materialmente relevante em magnitude,
com reduções de 34\% no MSE e de 24\% no MAE, e uma diferença de 16 pontos percentuais
no coeficiente de determinação da regressão de Mincer-Zarnowitz. A significância estatística
formal, avaliada pelo teste de Diebold-Mariano, fica na fronteira do nível convencional
de 5\% (p = 0,0503), possivelmente em razão do limitado número de dias sobrepostos
disponíveis para comparação; a análise dos intervalos de confiança bootstrap para o QLIKE
— a métrica teoricamente mais adequada — fornece evidência mais direta da superioridade
das dollar bars, com intervalos não sobrepostos a 95\% de confiança.
